% CV scientifique 1 page (FR) — El Mehdi Haress
% Cible : ILS Data Scientist (jumeau numérique, maintenance prévisionnelle) — KNDS France, Versailles
% Compilation : xelatex cv-fr-knds.tex
\documentclass[10pt,a4paper]{article}

\usepackage{fontspec}
\usepackage[french]{babel}
\usepackage[T1]{fontenc}
\usepackage[
  top=0.85cm, bottom=0.85cm, left=1.4cm, right=1.4cm
]{geometry}
\usepackage{amsmath}
\usepackage{amssymb}
\usepackage{xcolor}
\usepackage{tikz}
\usepackage{enumitem}
\usepackage{graphicx}
\usepackage{titlesec}
\usepackage[hidelinks]{hyperref}
\usepackage{fontawesome5}

\setmainfont{texgyrepagella}[
  Extension      = .otf,
  UprightFont    = *-regular,
  BoldFont       = *-bold,
  ItalicFont     = *-italic,
  BoldItalicFont = *-bolditalic,
]
\newfontfamily\headingfont{texgyreheros}[
  Extension      = .otf,
  UprightFont    = *-regular,
  BoldFont       = *-bold,
  ItalicFont     = *-italic,
  BoldItalicFont = *-bolditalic,
]

\definecolor{accent}{RGB}{28,58,94}
\definecolor{ink}{RGB}{26,26,26}
\definecolor{muted}{RGB}{95,105,120}
\definecolor{rulegray}{RGB}{200,206,214}
\color{ink}

\pagestyle{empty}
\setlength{\parindent}{0pt}
\linespread{1.00}
\frenchspacing

\titleformat{\section}
  {\headingfont\bfseries\fontsize{9.5}{11}\selectfont\color{accent}}
  {}{0em}{\MakeUppercase}[\vspace{-0.85ex}{\color{rulegray}\rule{\linewidth}{0.7pt}}]
\titlespacing*{\section}{0pt}{1.6ex plus .2ex}{0.9ex}

\newlength{\datecol}\setlength{\datecol}{0.155\textwidth}
\newlength{\gutter}\setlength{\gutter}{0.85em}
\newlength{\bodycol}
\setlength{\bodycol}{\textwidth}
\addtolength{\bodycol}{-\datecol}
\addtolength{\bodycol}{-\gutter}

\newcommand{\entry}[3]{%
  \par\vspace{0.55ex}\noindent
  \begin{minipage}[t]{\datecol}
    \vspace{0pt}\raggedright\footnotesize\color{muted}#1
  \end{minipage}%
  \hspace{\gutter}%
  \begin{minipage}[t]{\bodycol}
    \vspace{0pt}\raggedright
    {\headingfont\bfseries\footnotesize #2}\\[0.08ex]
    {\footnotesize\color{accent}#3}
  \end{minipage}%
  \par\vspace{0.12ex}
}

\newenvironment{bullets}
  {\begin{list}{\textcolor{accent}{\scriptsize$\blacktriangleright$}}{%
      \setlength{\leftmargin}{\datecol}
      \addtolength{\leftmargin}{\gutter}
      \addtolength{\leftmargin}{1.15em}
      \setlength{\rightmargin}{0pt}
      \setlength{\labelwidth}{0.85em}
      \setlength{\labelsep}{0.3em}
      \setlength{\itemsep}{0.08ex}
      \setlength{\parsep}{0pt}
      \setlength{\topsep}{0.2ex}
      \setlength{\partopsep}{0pt}
      \footnotesize}}
  {\end{list}}

\newcommand{\skillrow}[2]{%
  \par\noindent
  \begin{minipage}[t]{\datecol}
    \vspace{0pt}\raggedright{\headingfont\bfseries\footnotesize #1}
  \end{minipage}%
  \hspace{\gutter}%
  \begin{minipage}[t]{\bodycol}
    \vspace{0pt}\raggedright\footnotesize #2
  \end{minipage}%
  \par\vspace{0.38ex}
}

\newcommand{\cpp}{\mbox{C++}}

\newcommand{\photo}[1]{%
  \begin{tikzpicture}
    \clip[rounded corners=4pt] (0,0) rectangle (2.35,3.008);
    \node[anchor=south west, inner sep=0] at (0,0)
      {\includegraphics[width=2.35cm]{#1}};
  \end{tikzpicture}%
}

\begin{document}

\begin{minipage}[t]{0.78\textwidth}
\vspace{0pt}
{\headingfont\bfseries\fontsize{20}{23}\selectfont EL MEHDI HARESS}\\[0.6ex]
{\headingfont\fontsize{10}{12}\selectfont\color{accent}%
Docteur en mathématiques appliquées \textbar{} Ingénieur CentraleSupélec}\\[0.25ex]
{\small\color{muted}Modélisation \& simulation \textbar{} Data science \textbar{} Quantification d'incertitude \textbar{} Python}\\[1.0ex]
{\footnotesize
\faEnvelope[regular]~\href{mailto:el-mehdi.haress@inria.fr}{el-mehdi.haress@inria.fr}\quad
\faPhone*~+33 7 58 23 41 50\\[0.25ex]
\faLinkedin~\href{https://www.linkedin.com/in/el-mehdi-haress-67750b3aa/}{el-mehdi-haress}\quad
\faGlobe~\href{https://elmehdiharess.org}{elmehdiharess.org}\quad
\faGithub~\href{https://github.com/ElMehdiHaress}{ElMehdiHaress}\\[0.25ex]
\faMapMarker*~Sophia Antipolis \textemdash{} mobilité Île-de-France (Versailles)}
\end{minipage}\hfill
\begin{minipage}[t]{0.18\textwidth}
\vspace{0pt}\raggedleft
\photo{assets/photo.jpg}
\end{minipage}

\vspace{0.7ex}

\section{Profil}
\footnotesize
\textbf{5 ans} de modélisation de
\textbf{systèmes dynamiques complexes} : simulation numérique avec garanties de convergence,
\textbf{calibration de modèles sur données observées}, analyse de sensibilité et apprentissage robuste aux
événements rares (AISTATS). Huit publications, code \emph{open source}, applications web livrées de
bout en bout. Objectif : mettre cette rigueur au service d'un
\textbf{jumeau numérique} qui \textbf{couple modèles physiques, données d'essais et d'exploitation, et IA}
pour anticiper le comportement des systèmes et éclairer les décisions de maintenance.
\normalsize

\section{Compétences}
\skillrow{Modélisation}{Modèles dynamiques stochastiques (EDS, EDPS), modèles lagrangiens, systèmes
multi-échelles, schémas de simulation et analyse de convergence, Monte-Carlo.}
\skillrow{Données \& IA}{Estimation de paramètres sur observations discrètes, inférence statistique,
quantification d'incertitude, sensibilité ; apprentissage robuste aux événements rares ; bases en
réseaux de neurones.}
\skillrow{Programmation}{\textbf{Python} (NumPy, SciPy, pandas) \textemdash{} 8 ans ; MATLAB, R, Git,
\LaTeX{} ; JavaScript/TypeScript, React, Firebase ; bases en \cpp{}.}
\skillrow{Savoir-être}{Traduire une question métier en problème algorithmique ; travail avec
mathématiciens, ingénieurs et partenaires industriels ; rédaction et restitution orale (15+ exposés).}
\skillrow{Langues}{Français \textemdash{} langue maternelle. Anglais \textemdash{} courant. Japonais \textemdash{} notions.}

\section{Expérience}

\entry{09/2026 \textemdash{} auj.}{Ingénieur de recherche \textemdash{} modélisation stochastique appliquée}{Inria Sophia Antipolis \textemdash{} équipe CALISTO (Mireille Bossy)}
\begin{bullets}
  \item Modèles stochastiques lagrangiens (turbulence) et systèmes multi-échelles industriels : schémas de
  simulation, implémentation Python, validation numérique.
\end{bullets}

\entry{12/2024 \textemdash{} 08/2026}{Chercheur post-doctorant en analyse stochastique}{University of Leeds (Royaume-Uni) \textemdash{} K. Dareiotis et K. Lê}
\begin{bullets}
  \item \textbf{Sensibilité} de dynamiques singulières aux perturbations du modèle et des paramètres (\emph{SPA}, 2026).
  \item Co-organisation du workshop FRSSD 2026 (18 intervenants). TD de processus stochastiques (licence, 50 étudiants).
\end{bullets}

\entry{10/2021 \textemdash{} 11/2024}{Doctorat \textemdash{} simulation numérique de dynamiques stochastiques singulières}{Université Paris-Saclay / CentraleSupélec \textemdash{} L. Goudenège et A. Richard}
\begin{bullets}
  \item Schémas de simulation avec garanties de convergence là où les méthodes standard échouent
  (\emph{IMA J. Numerical Analysis}, \emph{SPA}).
  \item \textbf{Calibration conjointe de plusieurs paramètres} (dérive, volatilité, mémoire) à partir de
  données discrètes ; implémentation et validation numérique, 5 dépôts publics (Python, MATLAB).
  \item Enseignement à CentraleSupélec ($\sim$60 étudiants/an) : probabilités, EDP, finance stochastique (M2).
\end{bullets}

\entry{03/2020 \textemdash{} 09/2021}{Stage de recherche \textemdash{} apprentissage automatique robuste}{Université d'Osaka (Japon) \textemdash{} M. Holland}
\begin{bullets}
  \item Algorithmes d'apprentissage sensibles au risque, robustes aux pertes extrêmes et aux queues lourdes ;
  2 articles \emph{AISTATS} avec implémentation et campagnes numériques.
\end{bullets}

\section{Formation}
\skillrow{2021--2024}{Doctorat en mathématiques appliquées \textemdash{} Université Paris-Saclay.}
\skillrow{2017--2021}{Ingénieur, majeure mathématiques appliquées \textemdash{} CentraleSupélec. M2 Mathématiques de l'aléatoire (2020--2021).}
\skillrow{2015--2017}{CPGE MP \textemdash{} Lycée Pierre de Fermat, Toulouse.}

\section{Publications sélectionnées}
{\scriptsize\color{muted}8 publications ; liste complète sur \href{https://elmehdiharess.org}{elmehdiharess.org}. Ordre alphabétique en probabilités.\par}
\vspace{0.25ex}
\begin{itemize}[leftmargin=1.15em, itemsep=0.12ex, topsep=0pt, parsep=0pt, label={\textcolor{accent}{\scriptsize$\blacktriangleright$}}]
  \scriptsize
  \item M. Holland, \textbf{E.M. Haress}. \emph{Spectral risk-based learning using unbounded losses}. AISTATS, 2022.
  \item \textbf{E.M. Haress}, A. Richard. \emph{Estimation of several parameters in discretely-observed SDEs
  with additive fractional noise}. Statistical Inference for Stochastic Processes, 2024.
  \item K. Dareiotis, \textbf{E.M. Haress}, K. Lê. \emph{Uniform pathwise stability of additive singular SDEs
  driven by fractional Brownian motion}. SPA, 2026.
\end{itemize}

\section{Compléments}
\skillrow{Projets}{Applications web conçues et mises en production en autodidacte : \textbf{Nerduncover}
(jeu multijoueur), \href{https://hanabellaspa.pro}{hanabellaspa.pro} (site client et back-office : comptes,
rôles, rapports, export CSV). React, TypeScript, Firebase.}

\end{document}
